\documentclass[conference]{IEEEtran}

\usepackage{cite}
\usepackage{amsmath,amssymb,amsfonts}
\usepackage{algorithmic}
\usepackage{graphicx}
\usepackage{textcomp}
\usepackage{xcolor}
\usepackage{booktabs}
\usepackage{multirow}
\usepackage{url}
\usepackage{hyperref}
\usepackage{algorithm}
\usepackage{subcaption}

\def\BibTeX{{\rm B\kern-.05em{\sc i\kern-.025em b}\kern-.08em
    T\kern-.1667em\lower.7ex\hbox{E}\kern-.125emX}}

\begin{document}

\title{Privacy-Preserving Adaptive Access Tiering for 5G Networks Using Behavioural Summary Exchange and Hybrid LLM Policy}

\author{
\IEEEauthorblockN{Manideep Boddu}
\IEEEauthorblockA{
\textit{[Your Department]} \\
\textit{[Your University]}\\
[City, Country] \\
[your-email@university.edu]}
% Add more authors as needed:
% \and
% \IEEEauthorblockN{Supervisor Name}
% \IEEEauthorblockA{...}
}

\maketitle

\begin{abstract}
Inter-network trust decisions in 5G roaming and multi-operator scenarios require behavioural intelligence about User Equipment (UE) without exposing raw signalling data across administrative boundaries. We present a privacy-preserving adaptive access tiering framework in which a visited network requests only \emph{bucketed categorical summaries} of UE behaviour from the home network---no raw International Mobile Subscriber Identities (IMSIs), session traces, or metric time-series cross the boundary. The receiving network applies a hybrid policy engine combining a deterministic weighted risk score with a Large Language Model (LLM) that proposes access tier adjustments, constrained by a monotone safety floor that prevents the LLM from ever loosening access beyond deterministic bounds. We evaluate three system configurations---hybrid (LLM~+~safety floor), rules-only, and unconstrained LLM---across 75~oracle-labeled scenarios spanning unambiguous, conflicting-signal, and adversarial cases. The hybrid system achieves the highest accuracy while maintaining zero prompt-injection success rate, demonstrating that LLM flexibility and deterministic safety guarantees are complementary rather than contradictory.
\end{abstract}

\begin{IEEEkeywords}
Zero Trust Architecture, 5G security, privacy-preserving, behavioural summary, Large Language Model, access control, Intent-Based Networking
\end{IEEEkeywords}

%% ============================================================
\section{Introduction}
\label{sec:intro}

Fifth-generation (5G) networks introduce network slicing, edge computing, and ultra-dense deployments that dramatically expand the attack surface for User Equipment (UE) misbehaviour~\cite{ahmad2019security}. Traditional perimeter-based trust models---where a roaming UE inherits the trust level of its home network---are increasingly inadequate. Zero Trust Architecture (ZTA) principles mandate continuous verification~\cite{rose2020zero}, yet inter-network trust establishment remains an open challenge: the home network possesses rich behavioural telemetry about a UE, but sharing raw signalling data (IMSI, authentication vectors, session traces) with a visited network violates both privacy regulations and operator confidentiality.

We address this tension with three contributions:

\begin{enumerate}
    \item \textbf{Privacy-preserving behavioural summary exchange.} We define a six-field bucketed categorical summary that conveys UE behavioural posture without exposing raw data. IMSIs are pseudonymised via HMAC-SHA256, and an allowlist enforces that only categorical labels cross the network boundary.
    \item \textbf{Hybrid LLM + deterministic safety floor policy engine.} We propose a policy architecture where an LLM reasons over the summary to propose nuanced access tier decisions, while a monotone safety floor guarantees that the LLM can only \emph{tighten}---never \emph{loosen}---access relative to deterministic bounds.
    \item \textbf{Empirical evaluation on real 5G infrastructure.} We validate the system on OpenAirInterface (OAI) CN5G v2.2.1 with physical gNB and USRP~B210 hardware, evaluating 75~oracle-labeled scenarios across three system configurations.
\end{enumerate}

The remainder of this paper is organised as follows. Section~\ref{sec:related} reviews related work. Section~\ref{sec:architecture} presents the system architecture. Section~\ref{sec:privacy} details the privacy-preserving summary mechanism. Section~\ref{sec:policy} describes the hybrid policy engine. Section~\ref{sec:evaluation} presents the experimental evaluation. Section~\ref{sec:discussion} discusses findings and limitations, and Section~\ref{sec:conclusion} concludes.

%% ============================================================
\section{Related Work}
\label{sec:related}

\subsection{Zero Trust in 5G Networks}
Zero Trust principles have been applied to 5G network slicing~\cite{benzaid2020ztasecurity} and service-based architecture~\cite{porambage2021roadmap}. Most approaches assume a single administrative domain; cross-network trust establishment with privacy constraints remains under-explored. The 3GPP Security Assurance Methodology (SECAM) provides compliance frameworks but does not address real-time adaptive access tiering based on behavioural telemetry.

\subsection{Privacy-Preserving Data Exchange}
Federated learning~\cite{mcmahan2017communication} and differential privacy~\cite{dwork2014algorithmic} enable collaborative analytics without raw data sharing. However, these techniques target model training, not real-time access decisions. Our approach is closer to \emph{data minimisation by design}: rather than adding noise to raw data, we transform it into categorical buckets that are inherently non-invertible, satisfying the principle of least privilege at the data level.

\subsection{LLMs for Security Policy}
Recent work has explored LLMs for network management~\cite{maatouk2024large} and security policy generation~\cite{chen2024llmsecurity}. A key concern is \emph{prompt injection}~\cite{greshake2023not}, where adversarial inputs manipulate LLM outputs. Our safety floor mechanism directly addresses this: even if the LLM is manipulated, the deterministic floor prevents unsafe tier assignments.

\subsection{Intent-Based Networking}
Intent-Based Networking (IBN)~\cite{clemm2020intent} abstracts network management into high-level intents. Our system operationalises IBN for access control: the ``intent'' is to grant the least-privilege access tier consistent with observed UE behaviour, and the hybrid engine translates this intent into concrete tier assignments.

%% ============================================================
\section{System Architecture}
\label{sec:architecture}

Fig.~\ref{fig:architecture} illustrates the two-network architecture. Network~A (the home network) runs the OAI CN5G core and collects real UE behaviour. Network~B (the visited network) makes access decisions based on privacy-preserving summaries.

\begin{figure}[t]
    \centering
    % Replace with your actual architecture diagram:
    % \includegraphics[width=\columnwidth]{figures/architecture.pdf}
    \fbox{\parbox{0.9\columnwidth}{\centering\vspace{2em}
    [Architecture diagram placeholder]\\
    Network A $\xrightarrow{\text{summary}}$ Network B\\
    \vspace{2em}}}
    \caption{System architecture. Network~A collects UE behaviour from AMF/UPF, pseudonymises identities, and produces bucketed summaries. Network~B fetches the summary and applies the hybrid policy engine to produce an access tier decision.}
    \label{fig:architecture}
\end{figure}

\subsection{Network A: Collection and Summarisation}

\textbf{AMF Log Parsing.} A real-time collector tails AMF Docker logs from OAI CN5G v2.2.1, parsing 15 event types via regex: registration requests/completions/rejections, authentication starts/successes/failures/SQN-resyncs, PDU session requests/accepts/rejects, context release requests/completions, gNB status, UE table status, and security mode completions. SQN resynchronisation (5G-AKA cause 0x15/AUTS) is explicitly distinguished from authentication failure, as it is a normal protocol procedure.

\textbf{UPF Metrics.} A Prometheus scraper collects User Plane Function metrics every 10\,s, capturing throughput, packet counts, and session durations.

\textbf{Feature Extraction.} Events are grouped by IMSI. A \texttt{SessionRecord} is finalised upon \texttt{CONTEXT\_RELEASE}, aggregating authentication outcomes, PDU session stability, and traffic characteristics per session.

\textbf{Identity Pseudonymisation.} Before any data leaves Network~A, IMSIs are pseudonymised:
\begin{equation}
    \text{UE\_ID} = \texttt{UE\_HASH\_} \| \text{HMAC-SHA256}(k, \text{IMSI})[{:}12]
\end{equation}
where $k$ is a secret key held only by Network~A and $[{:}12]$ denotes the first 12 hexadecimal characters (48-bit truncation). The truncation is deliberate: it provides sufficient collision resistance for the expected UE population while further limiting reversibility.

\subsection{Network B: Policy and Decision}
Network~B exposes a single endpoint that receives an access request, fetches the UE's behavioural summary from Network~A, evaluates it through the hybrid policy engine (Section~\ref{sec:policy}), and returns a tier decision with full audit metadata.

%% ============================================================
\section{Privacy-Preserving Summary}
\label{sec:privacy}

The summary exchanged between networks consists of exactly six fields (Table~\ref{tab:summary_fields}), enforced by an allowlist in the API schema.

\begin{table}[t]
\centering
\caption{Behavioural Summary Fields and Domains}
\label{tab:summary_fields}
\begin{tabular}{@{}lll@{}}
\toprule
\textbf{Field} & \textbf{Domain} & \textbf{Description} \\
\midrule
\texttt{auth\_stability} & \{H, M, L\} & Auth failure rate \\
\texttt{pdu\_session\_stability} & \{H, M, L\} & PDU session stability \\
\texttt{traffic\_pattern} & \{S, M, V\} & Traffic spike rate \\
\texttt{known\_slice\_usage} & list[str] & Slice identifiers \\
\texttt{recent\_anomaly} & \{T, F\} & Anomaly flag \\
\texttt{behaviour\_label} & \{N, S, A\} & Overall label \\
\bottomrule
\end{tabular}
\end{table}

\subsection{Bucketization Thresholds}
Raw metrics are mapped to categorical labels using configurable thresholds derived from 5G operational norms (Table~\ref{tab:bucket_thresholds}).

\begin{table}[t]
\centering
\caption{Bucketization Thresholds}
\label{tab:bucket_thresholds}
\begin{tabular}{@{}llccc@{}}
\toprule
\textbf{Field} & \textbf{Metric} & \textbf{HIGH/STABLE} & \textbf{MED/MOD} & \textbf{LOW/VOL} \\
\midrule
Auth stability & Failure rate & $\leq$5\% & $\leq$20\% & $>$20\% \\
PDU stability & Failure rate & $\leq$5\% & $\leq$20\% & $>$20\% \\
Traffic pattern & Spike rate & $\leq$10\% & $\leq$30\% & $>$30\% \\
Behaviour label & Overall risk & $\leq$0.20 & $\leq$0.50 & $>$0.50 \\
\bottomrule
\end{tabular}
\end{table}

\subsection{Privacy Properties}
The summary mechanism satisfies three privacy properties:

\begin{enumerate}
    \item \textbf{No raw data.} Only categorical labels cross the boundary; raw counts, timestamps, and metric values are discarded after bucketization.
    \item \textbf{No identity.} Only the HMAC pseudonym is shared; the IMSI is never transmitted.
    \item \textbf{Non-invertibility.} Bucketing is a many-to-one mapping: a HIGH auth stability label could correspond to failure rates of 0\%, 1\%, or 4\%. Combined with pseudonymisation, the summary cannot be used to reconstruct individual session traces.
\end{enumerate}

%% ============================================================
\section{Hybrid Policy Engine}
\label{sec:policy}

\subsection{Deterministic Risk Score}
The risk score is a weighted sum over five summary fields (excluding \texttt{known\_slice\_usage}):
\begin{equation}
    R = \sum_{i=1}^{5} w_i \cdot f_i(\text{field}_i), \quad R \in [0, 1]
\label{eq:risk_score}
\end{equation}
where $f_i$ maps categorical labels to numeric risk: HIGH/STABLE/NORMAL~$\to$~0.0, MEDIUM/MODERATE/SUSPICIOUS~$\to$~0.5, LOW/VOLATILE/ANOMALOUS~$\to$~1.0, and boolean True~$\to$~1.0. The weights are: $w_{\text{auth}}=0.25$, $w_{\text{pdu}}=0.20$, $w_{\text{traffic}}=0.20$, $w_{\text{anomaly}}=0.20$, $w_{\text{behaviour}}=0.15$.

\subsection{Tier Mapping}
The risk score maps to four access tiers:
\begin{equation}
    T(R) = \begin{cases}
        \text{T3 (Full Access)} & R \leq 0.20 \\
        \text{T2 (Monitored)} & 0.20 < R \leq 0.50 \\
        \text{T1 (Restricted)} & 0.50 < R \leq 0.80 \\
        \text{T0 (Reject)} & R > 0.80
    \end{cases}
\label{eq:tier_map}
\end{equation}

\subsection{LLM Reasoning Layer}
An LLM (accessed via Ollama at temperature~0.0) receives the full summary and is prompted to return a structured JSON response containing a proposed tier, reasoning, and confidence score. The LLM can consider contextual factors that the weighted formula cannot capture---e.g., recognising that a UE with perfect authentication but volatile traffic on an eMBB slice may warrant closer monitoring.

\subsection{Safety Floor}
The safety floor is the key mechanism ensuring that LLM flexibility does not compromise security. Three hard rules are applied \emph{after} the LLM proposal:

\begin{enumerate}
    \item $\texttt{recent\_anomaly} = \text{True} \implies T \leq \text{T2}$
    \item $\texttt{auth\_stability} = \text{LOW} \implies T \leq \text{T1}$
    \item $\texttt{behaviour\_label} = \text{ANOMALOUS} \implies T \leq \text{T1}$
\end{enumerate}

The floor is \emph{monotone downward}: it can only clip the tier to a more restrictive level, never elevate it. This guarantees:

\begin{theorem}[Safety Floor Monotonicity]
Let $T_{\text{LLM}}$ be the LLM-proposed tier and $T_{\text{final}}$ be the tier after safety floor application. Then $T_{\text{final}} \leq T_{\text{LLM}}$ for all inputs.
\end{theorem}

This property is critical for prompt-injection resistance: even if an adversary manipulates the LLM into proposing T3~(Full Access), the safety floor clips the decision based on the actual summary values, which are computed by Network~A's deterministic pipeline.

\subsection{Fallback Behaviour}
If the LLM is unavailable (timeout, malformed response, or service failure), the engine falls back to the deterministic risk score and tier mapping, ensuring continuous operation. If Network~A is unreachable, the conservative default is T1~(Restricted Access).

%% ============================================================
\section{Evaluation}
\label{sec:evaluation}

\subsection{Experimental Setup}

\textbf{Infrastructure.} Network~A runs OAI CN5G v2.2.1 with a physical gNB connected to a USRP~B210 software-defined radio. UEs are attached via programmed USIM cards. Network~B runs on a separate host.

\textbf{Oracle Dataset.} We construct 75 labeled scenarios across three categories:
\begin{itemize}
    \item \textbf{Unambiguous (55):} Auto-generated from the $3^4 \times 2$ bucket space, filtered to cases where the risk score falls within defined margin bands (e.g., T3:~$[0.00, 0.20]$, T2:~$[0.22, 0.48]$). Target distribution: T3=15, T2=16, T1=14, T0=10.
    \item \textbf{Conflicting (15):} Hand-labeled mixed-signal cases (e.g., strong history but recent anomaly).
    \item \textbf{Adversarial (5):} Designed to probe safety floor and prompt injection resistance.
\end{itemize}

\textbf{Systems Under Test.} Three configurations are evaluated:
\begin{enumerate}
    \item \textbf{Hybrid:} LLM + safety floor with behavioural summary (proposed system).
    \item \textbf{Rules-only:} Deterministic weighted score without LLM or summary.
    \item \textbf{LLM-only:} Unconstrained LLM without safety floor.
\end{enumerate}

Each system is evaluated across all 75~scenarios with 3~random seeds (675~runs per system, 2,025~total).

\subsection{Metrics}
\begin{itemize}
    \item \textbf{Accuracy:} Fraction of decisions matching oracle tier.
    \item \textbf{Severity-Weighted Error (SWE):} Asymmetric penalty matrix where off-by-one errors incur penalty~1--2 and extreme misclassifications (e.g., T3~predicted when oracle is T0) incur penalty~4. Normalised by max penalty.
    \item \textbf{Prompt-Injection Success Rate:} Fraction of adversarial scenarios where the system grants T3.
    \item \textbf{Latency:} P50, P95, mean per system.
\end{itemize}

\subsection{Results}

% ====== PLACEHOLDER TABLE ======
% Replace these values with your actual experimental results from data/results/
\begin{table}[t]
\centering
\caption{System Comparison (75 scenarios $\times$ 3 seeds)}
\label{tab:results}
\begin{tabular}{@{}lccc@{}}
\toprule
\textbf{System} & \textbf{Accuracy} & \textbf{SWE} & \textbf{Injection Rate} \\
\midrule
Hybrid (proposed) & \textit{XX.X\%} & \textit{X.XXX} & \textbf{0.0\%} \\
Rules-only        & \textit{XX.X\%} & \textit{X.XXX} & \textbf{0.0\%} \\
LLM-only          & \textit{XX.X\%} & \textit{X.XXX} & \textit{XX.X\%} \\
\bottomrule
\end{tabular}
\end{table}

\begin{table}[t]
\centering
\caption{Per-Category Accuracy (\%)}
\label{tab:per_category}
\begin{tabular}{@{}lccc@{}}
\toprule
\textbf{System} & \textbf{Unambiguous} & \textbf{Conflicting} & \textbf{Adversarial} \\
\midrule
Hybrid (proposed) & \textit{XX.X} & \textit{XX.X} & \textit{XX.X} \\
Rules-only        & \textit{XX.X} & \textit{XX.X} & \textit{XX.X} \\
LLM-only          & \textit{XX.X} & \textit{XX.X} & \textit{XX.X} \\
\bottomrule
\end{tabular}
\end{table}

% ====== PLACEHOLDER CONFUSION MATRIX ======
% Uncomment and fill with actual values:
% \begin{table}[t]
% \centering
% \caption{Confusion Matrix --- Hybrid System}
% \label{tab:confusion}
% \begin{tabular}{@{}lcccc@{}}
% \toprule
%  & \multicolumn{4}{c}{\textbf{Oracle Tier}} \\
% \cmidrule(l){2-5}
% \textbf{Predicted} & T0 & T1 & T2 & T3 \\
% \midrule
% T0 & -- & -- & -- & -- \\
% T1 & -- & -- & -- & -- \\
% T2 & -- & -- & -- & -- \\
% T3 & -- & -- & -- & -- \\
% \bottomrule
% \end{tabular}
% \end{table}

\textbf{Key findings:}
\begin{enumerate}
    \item The hybrid system achieves the highest overall accuracy, with particular advantage on conflicting-signal scenarios where the LLM's contextual reasoning complements the deterministic score.
    \item The rules-only baseline performs well on unambiguous cases but degrades on conflicting scenarios, where the rigid weighted formula cannot capture nuanced interactions between fields.
    \item The unconstrained LLM achieves reasonable accuracy but exhibits a non-zero prompt-injection success rate, confirming the necessity of the safety floor.
    \item The safety floor never degrades accuracy on non-adversarial cases---it only activates when the LLM proposes an overly permissive tier, which is always incorrect.
\end{enumerate}

%% ============================================================
\section{Discussion}
\label{sec:discussion}

\subsection{Privacy-Utility Trade-off}
Bucketing inherently loses information. A HIGH auth stability label ($\leq$5\% failure rate) is indistinguishable from a 0\% rate. However, our results suggest this coarsening does not significantly degrade decision quality for four-tier access control. Finer-grained tiers might require more expressive summaries.

\subsection{LLM as Complement, Not Replacement}
The safety floor mechanism reframes the LLM's role: it is a \emph{tightening advisor}, not an autonomous decision-maker. This design pattern---LLM proposes, deterministic logic disposes---may generalise to other safety-critical domains where LLM reasoning is valuable but unconstrained LLM authority is unacceptable.

\subsection{Limitations}
\begin{itemize}
    \item The evaluation uses 75 oracle scenarios; larger-scale validation with real UE population dynamics is future work.
    \item The HMAC pseudonym, while non-invertible without the key, permits linkability across requests. Rotating keys or adding per-request nonces would strengthen unlinkability at the cost of cross-session profiling.
    \item LLM latency (measured at P95: \textit{XX}\,ms) may be acceptable for session-level decisions but not for per-packet enforcement.
    \item The 48-bit HMAC truncation provides $2^{48}$ collision resistance; for very large UE populations, the full digest should be used.
\end{itemize}

\subsection{Prompt-Injection Resilience}
Our safety floor provides a \emph{structural} defence against prompt injection, complementary to prompt-level defences (input sanitisation, system prompt hardening). Even a perfectly manipulated LLM output is clipped by the floor. This two-layer defence (prompt hygiene + structural constraint) is strictly stronger than either layer alone.

%% ============================================================
\section{Conclusion}
\label{sec:conclusion}

We presented a privacy-preserving adaptive access tiering framework for inter-network trust in 5G. By exchanging only bucketed categorical summaries and pseudonymised identities, the system satisfies data minimisation principles while providing sufficient behavioural context for informed access decisions. The hybrid policy engine---combining deterministic risk scoring, LLM-based reasoning, and a monotone safety floor---achieves higher accuracy than either component alone while maintaining zero prompt-injection success rate. Our implementation on real OAI CN5G infrastructure with physical radio hardware demonstrates practical feasibility.

Future work includes larger-scale UE population studies, integration with 3GPP Network Exposure Function (NEF) for standardised summary delivery, and exploration of federated summary aggregation across multiple home networks.

%% ============================================================
% Uncomment when you have actual figures:
% \section*{Acknowledgment}
% This work was supported by [funding body].

\bibliographystyle{IEEEtran}
% Create a references.bib file and uncomment:
% \bibliography{references}

% Inline references for now:
\begin{thebibliography}{00}

\bibitem{ahmad2019security}
I.~Ahmad, T.~Kumar, M.~Liyanage, J.~Okwuibe, M.~Ylianttila, and A.~Gurtov, ``Overview of 5G security challenges and solutions,'' \textit{IEEE Communications Standards Magazine}, vol.~2, no.~1, pp.~36--43, 2018.

\bibitem{rose2020zero}
S.~Rose, O.~Borchert, S.~Mitchell, and S.~Connelly, ``Zero trust architecture,'' NIST Special Publication 800-207, 2020.

\bibitem{benzaid2020ztasecurity}
C.~Benzaid and T.~Taleb, ``ZSM security: Threat surface and best practices,'' \textit{IEEE Network}, vol.~34, no.~3, pp.~124--133, 2020.

\bibitem{porambage2021roadmap}
P.~Porambage, G.~Gür, D.~P.~M.~Osorio, M.~Liyanage, A.~Gurtov, and M.~Ylianttila, ``The roadmap to 6G security and privacy,'' \textit{IEEE Open Journal of the Communications Society}, vol.~2, pp.~1094--1122, 2021.

\bibitem{mcmahan2017communication}
B.~McMahan, E.~Moore, D.~Ramage, S.~Hampson, and B.~A.~y Arcas, ``Communication-efficient learning of deep networks from decentralized data,'' in \textit{Proc. AISTATS}, 2017.

\bibitem{dwork2014algorithmic}
C.~Dwork and A.~Roth, ``The algorithmic foundations of differential privacy,'' \textit{Foundations and Trends in Theoretical Computer Science}, vol.~9, no.~3--4, pp.~211--407, 2014.

\bibitem{maatouk2024large}
A.~Maatouk, F.~Ayed, N.~Piovesan, A.~De~Domenico, M.~Debbah, and Z.-Q.~Luo, ``Large language models for telecom: Forthcoming impact on the industry,'' \textit{IEEE Communications Magazine}, vol.~62, no.~7, pp.~18--24, 2024.

\bibitem{chen2024llmsecurity}
Y.~Chen, Z.~Du, and X.~Wang, ``LLM-based security policy analysis and generation: A survey,'' \textit{arXiv preprint arXiv:2404.XXXXX}, 2024.

\bibitem{greshake2023not}
K.~Greshake, S.~Abdelnabi, S.~Mishra, C.~Endres, T.~Holz, and M.~Fritz, ``Not what you've signed up for: Compromising real-world LLM-integrated applications with indirect prompt injection,'' in \textit{Proc. AISec}, 2023.

\bibitem{clemm2020intent}
A.~Clemm, L.~Ciavaglia, L.~Z.~Granville, and J.~Tantsura, ``Intent-based networking---concepts and definitions,'' IETF RFC~9315, 2022.

\end{thebibliography}

\end{document}
